% problem-000435 8 溴化汞氨配合物
\section*{8. 溴化汞氨配合物}
从A到D的结构中， $\operatorname { H g } ( \operatorname { I I } )$ 与N的结合随N上所连氢原⼦的数⽬⽽变化，N均成4个键， $\mathrm { N { - } H g { - } N }$ 键⻆为180°。A中，Br−作简单⽴⽅堆积，两个⽴⽅体共⽤的⾯中⼼存在⼀个 $\operatorname { H g } ( \operatorname { I I } )$ $\mathrm { N H _ { 3 } }$ 位于⽴⽅体的体$\therefore E _ { \circ }$ B 中，Hg(II)与氨基形成⼀维链。C 中存在 $\operatorname { H g } ( \operatorname { I I } )$ 与亚氨基形成的按六边形扩展的⼆维结构，Br−位于六边形中⼼。D中， $\mathrm { ( H g _ { 2 } N ) } ^ { . }$ +形成具有类似 $\mathrm { S i O _ { 2 } }$ 的三维结构。

\subsection*{8-1}
写出⽣成C和D的⽅程式。

\subsection*{8-2}
画出A的结构⽰意图 $\mathrm { ( N H _ { 3 } }$ 以N表⽰）。

\subsection*{8-3}
画出 B 中 $\operatorname { H g } ( \operatorname { I I } )$ 与氨基（以N表⽰）形成的⼀维链式结构⽰意图。

\subsection*{8-4}
画出C中⼆维层的结构⽰意图，写出其组成式。层间还存在哪些物种？给出其⽐例。

\subsection*{8-5}
画出 D 中 $\mathrm { ( H g _ { 2 } N ) }$ +的三维结构⽰意图(⽰出 $\mathrm { H g }$ 与N的连接⽅式即可)。

\subsection*{8-6}
令⼈惊奇的是，组成为 HgNH_{2}F 的化合物并⾮与 $\mathrm { { H g N H _ { 2 } X ( X = C l , \Delta \mathrm { { B r } } } }$ ，I)同构，⽽与D相似，存在三维结构的 $\mathrm { ( H g 2 N ) _ { \circ } ^ { + } }$ 写出表达 $\mathrm { H g N H _ { 2 } F }$ 结构特点的结构简式。
